\subsection{Khái niệm: workflow, job, step, runner và trigger}

GitHub Actions là nền tảng tích hợp liên tục (Continuous Integration, viết tắt là CI)
và phân phối liên tục (Continuous Delivery, viết tắt là CD) được tích hợp sẵn trong GitHub.
Mỗi khi có sự kiện như đẩy code (\verb|push|) hoặc mở pull request, GitHub tự động thực hiện
một chuỗi công việc kiểm tra, biên dịch và kiểm thử đã được mô tả sẵn \cite{ghactions}.

Đơn vị cấu hình nhỏ nhất là một tệp tin YAML đặt trong thư mục
\verb|.github/workflows| của repository, ví dụ \verb|ci-aspx-481.yml|.
Khi tệp tin này tồn tại trên nhánh mặc định, thẻ Actions trên GitHub sẽ liệt kê
lịch sử chạy, nhật ký từng bước và kết quả đạt hay không đạt của mỗi lần chạy.

\begin{description}
\item[Workflow] Một tiến trình tự động hoàn chỉnh, được định nghĩa bởi một tệp tin YAML.
Ví dụ: workflow \verb|CI ASPX 4.8.1| chạy mỗi khi có \verb|push| hoặc \verb|pull_request|.
Một repository có thể chứa nhiều workflow chạy độc lập cho các mục đích khác nhau.
\item[Job] Một công việc lớn trong workflow, gồm nhiều bước chạy nối tiếp trên cùng một máy ảo.
Các job mặc định chạy song song; dùng \verb|needs| để bắt job kiểm thử đợi job biên dịch xong.
Ví dụ: job \verb|build| biên dịch solution, job \verb|test| chạy kiểm thử đơn vị.
\item[Step] Một bước nhỏ trong job, hoặc chạy lệnh shell qua \verb|run| hoặc tái sử dụng hành động có sẵn qua \verb|uses|.
Ví dụ: bước khôi phục gói chạy lệnh \verb|nuget restore|, bước kiểm thử chạy \verb|vstest.console.exe|.
Mỗi bước có nhật ký riêng nên khi lỗi rất dễ xác định bước nào gây ra.
\item[Runner] Máy ảo hoặc máy thật thực thi job, do GitHub cung cấp hoặc nhóm tự quản lý.
Project ASPX 4.8.1 bắt buộc dùng runner \verb|windows-latest| vì cần MSBuild và .NET Framework.
Runner Ubuntu không biên dịch được project Web Forms dùng .NET Framework truyền thống.
\item[Trigger] Sự kiện khởi động workflow, khai báo trong khối \verb|on| của tệp tin YAML.
Hai trigger phổ biến nhất là \verb|push| và \verb|pull_request| kèm bộ lọc nhánh.
Ngoài ra còn trigger lịch định kỳ, trigger thủ công và trigger khi phát hành phiên bản mới.
\end{description}

Ba trigger được dùng nhiều nhất trong đồ án môn học là:
\begin{itemize}
\item \verb|push| lên các nhánh \verb|main| và \verb|develop| để kiểm tra mọi commit mới.
\item \verb|pull_request| nhắm vào nhánh \verb|main| để kiểm tra code trước khi hợp nhất.
\item Kích hoạt thủ công để chạy lại pipeline khi cần kiểm tra gấp mà không đẩy commit mới.
\end{itemize}

Hiểu rõ năm khái niệm trên giúp đọc được mọi tệp tin workflow: trigger quyết định khi nào chạy,
runner quyết định chạy ở đâu, còn job và step quyết định chạy cái gì theo thứ tự nào.
Phần tiếp theo trình bày một workflow CI hoàn chỉnh cho project ASPX 4.8.1 từ khâu
lấy code, khôi phục gói NuGet, biên dịch, kiểm thử đến lưu trữ kết quả.

\subsection{Workflow CI đầy đủ cho ASPX 4.8.1}

\subsubsection{Tệp workflow hoàn chỉnh}

Tệp tin \verb|ci-aspx-481.yml| dưới đây là workflow CI hoàn chỉnh cho một solution ASP.NET Web Forms
viết bằng C\# với mục tiêu .NET Framework 4.8.1. Workflow chạy trên runner \verb|windows-latest|,
gồm tám bước: lấy code, cài MSBuild, cài NuGet, dùng cache, khôi phục gói,
biên dịch chế độ Release, chạy kiểm thử và tải kết quả lên artifact.
Mỗi bước đều có chú thích tiếng Anh để tương thích với quy tắc chỉ dùng ký tự ASCII trong khối liệt kê code.

\begin{lstlisting}[language=yaml, caption={Workflow CI đầy đủ cho project ASPX 4.8.1}, label={lst:c42-ci}]
name: CI ASPX 4.8.1

# Trigger: run on push and pull requests to main branches.
on:
  push:
    branches: [ main, develop ]
  pull_request:
    branches: [ main ]

# Cancel older runs of the same branch or pull request.
concurrency:
  group: ci-${{ github.ref }}
  cancel-in-progress: true

jobs:
  build:
    # .NET Framework 4.8.1 requires a Windows runner.
    runs-on: windows-latest

    steps:
      # Step 1: check out the repository source code.
      - name: Checkout code
        uses: actions/checkout@v4

      # Step 2: install MSBuild on the runner PATH.
      - name: Setup MSBuild
        uses: microsoft/setup-msbuild@v2

      # Step 3: install NuGet CLI on the runner PATH.
      - name: Setup NuGet
        uses: nuget/setup-nuget@v2

      # Step 4: restore the NuGet cache to speed up restore.
      - name: Cache NuGet packages
        uses: actions/cache@v4
        with:
          path: ~/.nuget/packages
          key: nuget-${{ runner.os }}-${{ hashFiles('**/packages.config', '**/*.csproj') }}
          restore-keys: |
            nuget-${{ runner.os }}-

      # Step 5: restore all NuGet packages for the solution.
      - name: Restore NuGet packages
        run: nuget restore AspxDemo_NET481.sln

      # Step 6: build the solution in Release mode.
      - name: Build solution
        run: msbuild AspxDemo_NET481.sln /p:Configuration=Release /p:Platform="Any CPU" /p:DeployOnBuild=true /p:PublishDir=publish\

      # Step 7: run unit tests with VSTest.
      - name: Run tests
        run: vstest.console.exe Tests\bin\Release\AspxDemo.Tests.dll /Logger:trx

      # Step 8: upload the published site as an artifact.
      - name: Upload artifact
        uses: actions/upload-artifact@v4
        with:
          name: aspx-site
          path: publish/
          retention-days: 7
\end{lstlisting}

Workflow trên có thể sao chép vào thư mục \verb|.github/workflows| của repository demo
và chạy ngay sau khi đẩy lên GitHub, với điều kiện solution có tên \verb|AspxDemo_NET481.sln|
và project kiểm thử đã được biên dịch ra \verb|AspxDemo.Tests.dll|.

\subsubsection{Giải thích từng bước}
\begin{enumerate}
\item \textbf{Lấy code:} hành động \verb|actions/checkout@v4| tải mã nguồn tại đúng commit kích hoạt.
Không có bước này, các bước sau không có code để biên dịch và kiểm thử.
\item \textbf{Cài MSBuild:} hành động \verb|microsoft/setup-msbuild@v2| đưa \verb|MSBuild.exe| vào biến PATH.
MSBuild là công cụ duy nhất biên dịch đầy đủ project Web Forms mục tiêu .NET Framework 4.8.1.
\item \textbf{Cài NuGet:} hành động \verb|nuget/setup-nuget@v2| đưa \verb|nuget.exe| vào biến PATH.
Lệnh \verb|nuget restore| sau đó đọc \verb|packages.config| và tải các gói còn thiếu về máy chạy.
\item \textbf{Dùng cache:} hành động \verb|actions/cache@v4| lưu thư mục \verb|~/.nuget/packages| giữa các lần chạy.
Khóa cache gồm hệ điều hành và giá trị băm của tệp khóa gói nên khi không đổi thư viện thì bỏ qua tải lại.
\item \textbf{Khôi phục gói:} lệnh \verb|nuget restore AspxDemo_NET481.sln| tải phần gói còn thiếu sau khi đã dùng cache.
Bước này bắt buộc đứng trước biên dịch, nếu thiếu gói thì MSBuild sẽ báo lỗi ngay.
\item \textbf{Biên dịch:} lệnh \verb|msbuild AspxDemo_NET481.sln /p:Configuration=Release| biên dịch toàn bộ solution.
Mọi lỗi biên dịch làm workflow báo đỏ và chặn pull request cho đến khi được sửa.
\item \textbf{Kiểm thử:} lệnh \verb|vstest.console.exe| chạy các DLL kiểm thử và xuất nhật ký định dạng \verb|.trx|.
Chỉ cần một test không đạt, toàn bộ workflow được đánh dấu thất bại.
\item \textbf{Lưu artifact:} hành động \verb|actions/upload-artifact@v4| lưu gói triển khai trong 7 ngày.
Dùng tham số \verb|retention-days| để kiểm soát dung lượng lưu trữ của repository.
\end{enumerate}

Với tám bước trên, mọi pull request đều được biên dịch và kiểm thử tự động trước khi review.
Nhóm nên bật quy tắc yêu cầu workflow đạt mới được hợp nhất để nhánh chính luôn biên dịch được.
Phần tiếp theo trình bày cách quản lý secret và môi trường triển khai an toàn.

\subsection{Secrets và Environments}

Secret là biến môi trường được mã hóa, dùng để lưu thông tin nhạy cảm như chuỗi kết nối cơ sở dữ liệu,
mật khẩu tài khoản triển khai hoặc khóa ký gói. Giá trị secret không bao giờ hiển thị trong nhật ký:
GitHub tự động che thành \verb|***| khi secret vô tình được in ra. Khai báo secret tại thẻ Settings,
mục Secrets and variables, nút New repository secret; secret cấp môi trường thì khai báo trong từng environment.

\begin{itemize}
\item Secret cấp repository: mọi workflow trong repository đều đọc được.
Dùng cho khóa dùng chung như token đọc gói nội bộ của nhóm.
\item Secret cấp môi trường: chỉ workflow gắn đúng environment mới đọc được.
Dùng cho chuỗi kết nối, vì staging và production dùng cơ sở dữ liệu khác nhau.
\item Biến môi trường thường: giá trị hiển thị công khai trong nhật ký.
Chỉ dùng cho dữ liệu không nhạy cảm như tên thư mục đầu ra hoặc cờ bật chi tiết nhật ký.
\end{itemize}

Trong tệp tin YAML, đọc secret bằng cú pháp \verb|${{ secrets.TEN_SECRET }}|, ví dụ
\verb|${{ secrets.STAGING_CONN_STRING }}|. Cú pháp này do GitHub thay thế lúc chạy,
nên tuyệt đối không ghi giá trị thật vào code mà chỉ ghi tên secret.
Tài liệu chính thức tại \url{https://docs.github.com/actions} mô tả đầy đủ cú pháp và phạm vi áp dụng.

\begin{lstlisting}[language=yaml, caption={Đọc secret theo môi trường staging và production}, label={lst:c42-secrets}]
jobs:
  deploy-staging:
    # Deploy job bound to the staging environment.
    runs-on: windows-latest
    environment: staging
    steps:
      # Read the staging connection string from secrets.
      - name: Deploy to staging
        run: deploy.cmd ${{ secrets.STAGING_CONN_STRING }}
  deploy-production:
    # Production requires manual approval in settings.
    runs-on: windows-latest
    environment: production
    steps:
      # Read the production connection string from secrets.
      - name: Deploy to production
        run: deploy.cmd ${{ secrets.PROD_CONN_STRING }}
\end{lstlisting}

Environment gom secret, quy tắc bảo vệ và lịch sử triển khai của một môi trường.
Môi trường staging cho phép triển khai tự động mỗi khi hợp nhất vào nhánh develop để kiểm thử sớm.
Môi trường production yêu cầu ít nhất một người review phê duyệt trước khi triển khai,
đồng thời giới hạn chỉ nhánh chính mới được triển khai để tránh phát hành nhầm code thử nghiệm.

\begin{itemize}
\item Không bao giờ đẩy tệp tin chứa secret thật như bản \verb|Web.config| trên máy cá nhân lên repository.
\item Xoay vòng secret ngay khi thành viên rời nhóm hoặc nghi ngờ lộ, rồi cập nhật lại trên GitHub.
\item Dùng secret khác nhau cho staging và production để lộ môi trường thử nghiệm không ảnh hưởng hệ thống thật.
\end{itemize}

\subsection{Cache NuGet và artifacts}

Mỗi lần pipeline chạy, lệnh khôi phục gói phải tải hàng trăm MB thư viện từ Internet,
chiếm phần lớn thời gian và dễ thất bại khi mạng chậm. Cache lưu thư mục gói toàn cục
\verb|~/.nuget/packages| sau lần chạy đầu, các lần sau chỉ tải phần thay đổi nên thường
rút ngắn bước khôi phục từ vài phút xuống còn vài chục giây \cite{githubdocs}.

Khóa cache gồm tên hệ điều hành và giá trị băm của các tệp tin khóa gói như
\verb|packages.lock.json| hoặc \verb|packages.config|. Khi không đổi thư viện, khóa trùng khớp
và cache được dùng lại; khi thêm hoặc nâng cấp gói, khóa thay đổi nên cache cũ bị bỏ,
đồng thời khóa dự phòng bảo đảm vẫn tái sử dụng được phần lớn gói đã lưu.

\begin{lstlisting}[language=yaml, caption={Cache NuGet và tải kết quả kiểm thử lên artifact}, label={lst:c42-cache}]
steps:
  # Cache the global NuGet folder between runs.
  - name: Cache NuGet packages
    uses: actions/cache@v4
    with:
      path: ~/.nuget/packages
      key: nuget-${{ runner.os }}-${{ hashFiles('**/*.csproj', '**/packages.config') }}
      restore-keys: |
        nuget-${{ runner.os }}-

  # Packages missing from cache are downloaded by restore.
  - name: Restore packages
    run: nuget restore AspxDemo_NET481.sln

  # Keep the build output for download and deploy jobs.
  - name: Upload test results
    uses: actions/upload-artifact@v4
    with:
      name: test-results
      path: TestResults/*.trx
\end{lstlisting}

Artifact là nơi lưu kết quả đầu ra của pipeline: gói triển khai website, tệp tin nhật ký kiểm thử
định dạng \verb|.trx| hoặc bản dựng để job triển khai tải về dùng tiếp. Khác với cache vốn có thể
bị xóa bất cứ lúc nào, artifact được đảm bảo tồn tại đúng số ngày đã cấu hình qua \verb|retention-days|.
Nhờ artifact, thành viên review có thể tải đúng bản đã kiểm thử về chạy thử mà không cần biên dịch lại.

\begin{itemize}
\item Luôn đặt cache ngay trước bước khôi phục gói để bước này hưởng lợi trực tiếp.
\item Đưa tệp tin khóa gói vào repository để khóa cache ổn định giữa các lần chạy.
\item Chỉ tải artifact thật sự cần thiết và đặt thời gian lưu ngắn để tiết kiệm dung lượng.
\end{itemize}

\subsection{Runner Windows và Ubuntu cho .NET Framework 4.8.1}

Project ASP.NET Web Forms mục tiêu .NET Framework 4.8.1 chỉ biên dịch đầy đủ trên Windows,
vì bộ công cụ gồm MSBuild, trình biên dịch trang ASPX và thư viện tham chiếu .NET Framework
đều là thành phần gắn chặt với hệ điều hành Windows \cite{msmsbuild}. Nền tảng .NET Framework
khác với .NET hiện đại đa nền tảng: thư viện và công cụ của nó không có bản Linux \cite{msnetfx}.

Runner \verb|ubuntu-latest| khởi động nhanh và miễn phí nhiều phút hơn, nhưng không có MSBuild đầy đủ,
không có .NET Framework 4.8.1 và không chạy được \verb|vstest.console.exe|. Lệnh \verb|dotnet build|
trên Ubuntu chỉ biên dịch project SDK-style dùng .NET hiện đại, không hiểu project Web Forms
kiểu cũ với tệp tin \verb|packages.config| và quy trình biên dịch trang ASPX.

\begin{table}[H]
\centering
\caption{So sánh runner cho project ASPX 4.8.1}
\label{tab:c42-runner}
\begin{tabular}{p{3.4cm}p{4.2cm}p{4.2cm}}
\toprule
Tiêu chí & Runner Windows & Runner Ubuntu \\
\midrule
Hệ điều hành & Windows Server 2022 & Ubuntu 22.04 \\
MSBuild đầy đủ & Có sẵn qua \texttt{setup-msbuild} & Không có \\
Biên dịch .NET Framework 4.8.1 & Hỗ trợ đầy đủ & Không hỗ trợ \\
Trang ASPX Web Forms & Biên dịch được & Không biên dịch được \\
Kiểm thử VSTest & Chạy được & Không chạy được \\
Tốc độ và chi phí & Khởi động chậm hơn & Nhanh và rẻ hơn \\
Khuyến nghị & Bắt buộc cho ASPX 4.8.1 & Chỉ dùng cho .NET hiện đại \\
\bottomrule
\end{tabular}
\end{table}

Bảng \ref{tab:c42-runner} tóm tắt sự khác biệt: với project ASPX 4.8.1, runner Windows là lựa chọn bắt buộc,
dù tốn thời gian khởi động hơn. Hình \ref{fig:c42-pipeline} minh họa chuỗi công việc CI
chạy trên runner này từ lúc lấy code đến khi có artifact.

\begin{figure}[H]
\centering
\begin{tikzpicture}[>=Stealth, thick, box/.style={draw, rounded corners, fill=blue!10, minimum width=2.2cm, minimum height=0.9cm, align=center, font=\small}]
\node[box] (a) {Checkout};
\node[box, right=1.0cm of a] (b) {Restore};
\node[box, right=1.0cm of b] (c) {Build};
\node[box, right=1.0cm of c] (d) {Test};
\draw[->] (a) -- node[midway, above, font=\footnotesize] {push} (b);
\draw[->] (b) -- node[midway, above, font=\footnotesize] {cache} (c);
\draw[->] (c) -- node[midway, above, font=\footnotesize] {pass} (d);
\caption{Chuỗi công việc CI cho project ASPX trên runner Windows}
\label{fig:c42-pipeline}
\end{figure}

Tóm lại, GitHub Actions giúp nhóm ASPX tự động hóa toàn bộ vòng kiểm tra mỗi khi đẩy code,
với điều kiện chọn đúng runner Windows cho .NET Framework 4.8.1. Kết hợp workflow CI đạt,
secret quản lý tập trung và cache hợp lý, nhóm có thể hợp nhất pull request nhanh mà vẫn an toàn.